% =====================================================================
%  Выводы по лабораторной работе.
% =====================================================================

\labpartbreak{Выводы}

В ходе выполнения лабораторной работы была разработана методика оценки
защищённости сторонних программных систем, охватывающая все семь
направлений анализа из теоретических основ задания, и применена на
практике к реальному веб-приложению -- варианту лабораторной работы
№2.1 -- с позиции внешнего наблюдателя, без изначального доверия к
исходному коду. По результатам 26 автоматизированных проверок обнаружено
шесть замечаний (открытая документация API, отсутствие двух заголовков
транспортной безопасности, фактически неработающее перенаправление на
HTTPS, отсутствие ограничения частоты запросов на регистрации,
отсутствие серверного экранирования сохраняемого содержимого), для
каждого из которых предложена конкретная мера усовершенствования, а
также сформулирован отдельный набор рекомендаций по безопасной
эксплуатации приложения, не требующих изменения кода.

Семантически работа показала принципиальное отличие оценки защищённости
«снаружи» от анализа «изнутри», выполненного в лабораторной работе №2.1:
разработчик системы знает, что определённый маршрут документации
API должен быть недоступен в эксплуатации, потому что так задумано в
коде; внешний наблюдатель этого не знает и обнаруживает лишь
наблюдаемый факт -- маршрут отвечает или не отвечает, -- и именно это
несоответствие между тем, что «задумано», и тем, что «реально доступно
по сети», регулярно оказывается источником уязвимостей в эксплуатации
реальных систем, а не ошибки в самом алгоритме шифрования или
хеширования. Отдельно подтверждено, что часть свойств системы (SQL-
инъекции, устойчивость JWT, шифрование данных) успешно выдержала
проверку именно потому, что эти свойства были заложены разработчиком
осознанно в лабораторной работе №2.1, тогда как обнаруженные замечания
относятся преимущественно к конфигурации развёртывания и к
неполной симметрии между схожими по назначению конечными точками
(\texttt{login} и \texttt{register}), а не к отсутствию защитных
механизмов как таковых.

С прагматической точки зрения работа показала, что содержательный
аудит защищённости не требует непременно тяжёлых готовых сканеров --
воспроизводимый набор небольших целевых скриптов, реализующих те же
категории проверок (разведка ресурсов, инъекционное тестирование,
проверка заголовков, проверка сессии), дал конкретные, проверяемые и
привязанные к реальному HTTP-ответу замечания, которые можно
использовать как регрессионный набор проверок при последующих
изменениях приложения. Также показано, что не всякое найденное
отклонение является уязвимостью в строгом смысле: находка по XSS была
сознательно классифицирована как «не эксплуатируется текущим клиентом,
но является скрытым риском», а не как подтверждённая уязвимость,
поскольку выданная API полезная нагрузка была дополнительно проверена
против исходного кода конкретного клиента -- корректная оценка
защищённости требует различать эти две степени серьёзности, а не
приравнивать любое отклонение от эталонного поведения к атаке.
